\chapter{An\'alisis de antecedentes y aportaci\'on realizada}\label{analanteced}

\section{Contexto del sector}

El mercado de distribuci\'on alimentaria en Espa\~na presenta una combinaci\'on de gran
competencia entre cadenas, alta sensibilidad al precio y fuerte fragmentaci\'on de oferta.
Durante los \'ultimos a\~nos, la inflaci\'on en alimentaci\'on ha incrementado la necesidad de
herramientas que permitan comprar mejor sin aumentar el tiempo invertido.

Aunque existen comparadores de precios consolidados, la mayor parte se centra en mostrar
informaci\'on de coste por tienda, sin integrar de forma nativa el impacto log\'istico de
desplazarse a varios establecimientos.

\section{Estado del arte en comparaci\'on de precios}

La evoluci\'on observada puede resumirse en tres generaciones:

\begin{enumerate}
	\item \textbf{Directorios de precios:} listados est\'aticos o semiactualizados, con
		poca cobertura y baja frecuencia de refresco.
	\item \textbf{Comparadores multisupermercado:} soluciones como Soysuper, Findit u OCU
		Market, con mejor cat\'alogo y funcionalidades de lista.
	\item \textbf{Personalizaci\'on con IA:} tendencia emergente en la que se introducen
		recomendaciones contextuales y asistencia conversacional.
\end{enumerate}

En Espa\~na, el salto completo hacia la tercera generaci\'on todav\'ia es limitado,
especialmente en el segmento de cesta alimentaria diaria.

\section{Tecnolog\'ias relevantes para BargAIn}

\subsection{Geoespacial y rutas}

PostGIS permite resolver consultas de proximidad con precisi\'on y rendimiento, mientras que
OR-Tools facilita modelar el problema de ordenaci\'on de paradas como una variante del TSP/VRP
en instancias peque\~nas (m\'aximo 3--4 paradas en nuestro caso).

\subsection{Ingesta de datos de precio}

La combinaci\'on Scrapy + Playwright es adecuada para cat\'alogos din\'amicos de supermercados.
Sin embargo, por robustez de producto, BargAIn no depende de una \'unica fuente y combina:

\begin{itemize}
	\item scraping programado,
	\item aportaciones de crowdsourcing,
	\item actualizaci\'on manual de comercios verificados.
\end{itemize}

\subsection{OCR y asistente}

Para OCR se adopta Google Cloud Vision API por mejor rendimiento pr\'actico en tickets
y fotos reales. Para el asistente se integra Gemini mediante backend proxy, con
guardrails para limitar respuestas al dominio de la compra.

\section{Aportaci\'on diferencial de BargAIn}

La aportaci\'on del proyecto se concreta en cuatro ejes:

\begin{itemize}
	\item \textbf{Optimizaci\'on multicriterio real} (precio + distancia + tiempo) en lugar de
		comparaci\'on aislada de precios.
	\item \textbf{Inclusi\'on de comercio local} mediante portal business y actualizaci\'on de
		precios/promociones.
	\item \textbf{Digitalizaci\'on de lista} por OCR con validaci\'on de usuario.
	\item \textbf{Interacci\'on conversacional} para consultas complejas de compra.
\end{itemize}

Estas capacidades, integradas en una sola plataforma, posicionan BargAIn en un espacio
de mercado poco cubierto por herramientas actuales.
